\section{Firmware}
\label{sec:firmware}

This section summarises the firmware's current implementation state at
thesis level. The maintained, code-level reference is
\texttt{sw/architecture.md} in the firmware repository -- it is updated in
the same commit as any structural firmware change, so it is authoritative
where this section and it disagree. Failure modes and their mitigations are
not repeated here; see \autoref{sec:fmea} for that. State as of firmware
version 0.1 (\texttt{version.h}), branch \texttt{sw\_dev}, commit
\texttt{cb868d6}.

\FloatBarrier

\subsection{Platform and design pattern}

The firmware runs on an STM32G474VET6 (Arm Cortex-M4, \qty{170}{\mega\hertz}, \qty{512}{\kibi\byte} flash, \qty{128}{\kibi\byte} RAM), chosen primarily for its \ac{HRTIM} (\qty{217}{\pico\second} resolution) needed for the converter \ac{PWM}. The application is bare-metal, event-driven and single-threaded -- no real-time operating system. The main loop drains one event per iteration from a 64-entry circular queue (\texttt{event.c}):

\begin{table}[h!]
    \centering
    \begin{tabular}{ll}
        \rowcolor{gray} \textbf{Event} & \textbf{Trigger} \\
        \texttt{EVENT\_IIC\_RX / IIC\_TX} & \ac{I2C} \ac{DMA} transfer complete \\
        \texttt{EVENT\_IIC\_ERROR} & \ac{I2C} NACK / bus error \\
        \texttt{EVENT\_EEPROM\_READ / EEPROM\_TIMEOUT} & \ac{EEPROM} read complete / no response \\
        \texttt{EVENT\_SM\_STEP} & State-machine tick (\qty{20}{\milli\second}) \\
        \texttt{EVENT\_BMS\_TIMER} & \ac{BMS} polling timer (\qty{1}{\second}) \\
    \end{tabular}
    \caption{Firmware event types}
\end{table}

\subsection{State machine and operating modes}

\texttt{statemachine.c} steps every \qty{20}{\milli\second} and dispatches
between 11 modes, handling \ac{UI} navigation, output-enable logic and the
protection interlock:

\begin{table}[h!]
    \centering
    \begin{tabular}{lp{9cm}}
        \rowcolor{gray} \textbf{Mode} & \textbf{Description} \\
        IDLE & Startup / menu navigation \\
        60V\_OUT & Regulated \qty{60}{\volt} output, held on with the OUT button (\ReqRef{R:60V}) \\
        10A\_OUT & Regulated current output up to \qty{10}{\ampere}, held on with the OUT button (\ReqRef{R:BSPD}) \\
        RESISTANCE\_1A & Milliohmmeter, pulsed \qty{1}{\ampere} excitation \\
        RESISTANCE\_1mA & Ohmmeter, \qty{1}{\milli\ampere} excitation \\
        ISOMETER & Insulation resistance test, high-voltage flyback active \\
        VOLTMETER & Passive voltage measurement \\
        AMPMETER & Passive current measurement (secondary half-bridge shorted) \\
        SETTINGS & Settings menu (submodes: MENU, BMS, DISPLAY, CALIBRATION, ADC1, ADC2, ABOUT) \\
        SHUTDOWN & Controlled power-down sequence \\
        CHARGE & Battery charging -- auto-entered from IDLE, not user-selectable \\
    \end{tabular}
    \caption{Statemachine operating modes (\texttt{statemachine.h})}
\end{table}

Per-mode data that used to be hand-kept in lockstep across four separate
switch statements (control mode, relay/shunt \ac{GPIO} mask, enable
\ac{GPIO}, entry flags, \ac{ADC} trigger) is centralised into one table,
\texttt{mode\_table.c} -- adding a mode or changing its behaviour means
editing one row instead of four call sites.

\subsection{Control loop}

The control loop runs from the \ac{ADC}1 \ac{DMA}-complete \ac{ISR}, not
directly from \ac{HRTIM}: the \ac{HRTIM} timer period triggers an \ac{ADC}
conversion at \qty{1}{\mega\hertz} with a $\div20$ post-scaler
(\qty{50}{\kilo\hertz}), and the \ac{ISR} decimates by 2, landing the
control loop at \qty{25}{\kilo\hertz} (\texttt{CTRL\_FREQ}). Every
integral gain and startup ramp is scaled by this rate, so changing the
switching frequency or post-scaler changes the effective loop tuning.

The controller (\texttt{ctrl\_PID\_control.c}) is a generic P/I/D with
integrator-clamping anti-windup; the D term is compiled out, so every loop
runs as PI. Eight tunable loops share one registry,
\texttt{ctrl\_pid\_table[]} in \texttt{ctrl\_main.c}, which
\texttt{ctrl\_main\_init()} and the \ac{CLI}'s \texttt{setPI}/
\texttt{saveEEPROM}/\texttt{readEEPROM} commands all iterate:

\begin{table}[h!]
    \centering
    \begin{tabular}{p{3.5cm}p{9cm}}
        \rowcolor{gray} \textbf{Loop} & \textbf{Used for} \\
        Buck & Buck voltage (resistance modes) \\
        Boost & \qty{60}{\volt} output \\
        Current & \qty{10}{\ampere} output \\
        BoostIoutLimit & Output current limiter on the boost path \\
        ChargeCurrent & CC phase of charging \\
        ChargeVoltage & CV tail of charging \\
        HVVoltage & ISOMETER high-voltage flyback voltage \\
        HVCurrent & ISOMETER leakage-current limiter \\
    \end{tabular}
    \caption{Tunable PI loops}
\end{table}

All eight loops still run placeholder gains -- see \autoref{sec:fmea} for
why that matters and what backs it up in the meantime.

\paragraph{Charging.} CHARGE is auto-entered from IDLE when a
\qtyrange{15}{24}{\volt} charger is present, no \ac{BMS} fault is latched, the
pack is below the end voltage and no lockout from an earlier abort is active.
Closing the output relay onto a converter output that is not at the charger
voltage, and a current loop that started at its maximum-current duty, caused
large inrush currents, so the start is a sequence (\texttt{charge\_seq.c}):
\emph{precharge} (relay open, converter run as an open-loop boost to
\qty{95}{\percent} of the charger voltage, since the converter output node is
not measured), \emph{close} (relay commanded closed while the duty is stepped to
the zero-current value), \emph{CC ramp} (current reference from
\qty{0}{\ampere} to \qty{1}{\ampere} over \qty{10}{\second}) and finally the
constant-voltage tail, which regulates on the highest cell voltage with a stepped
reference. The charge ends once the pack is at end voltage, the current has
tapered below \qty{100}{\milli\ampere} and the cell spread is at most
\qty{50}{\milli\volt}. The sequence is host-tested, but not yet validated on
hardware; \autoref{sec:fmea} lists what is open (unmeasured converter output,
relay timing, untuned gains).

\subsection{Peripherals}

\begin{table}[h!]
    \centering
    \begin{tabular}{p{3cm}p{6cm}p{2.5cm}}
        \rowcolor{gray} \textbf{IC} & \textbf{Role} & \textbf{Interface} \\
        BQ76905 & 4-cell \ac{LFP} \ac{BMS}: monitoring, protection, passive balancing & \ac{I2C}4 (0x08) \\
        ADS131M04 & 4-ch precision $\Sigma\Delta$ \ac{ADC} for isolation current & \ac{SPI}3 \\
        LMG2100R044 $\times$2 & \ac{GaN} half-bridge drivers, primary converter & \ac{HRTIM} \ac{GPIO} \\
        1EDN7512B & Gate driver, high-voltage flyback primary switch & \ac{HRTIM} \ac{GPIO} \\
        W25N01GV & 1\,Gbit QSPI NAND flash -- menu icon store, logging not yet wired up & QUADSPI \\
        ZD24C02B & \qty{2}{\kilo bit} \ac{I2C} \ac{EEPROM}, calibration/ID data & \ac{I2C}4 (0x50) \\
        TCA9554 & \ac{I2C} \ac{IO} expander, hardware ID reading & \ac{I2C}4 (0x20) \\
        LP5817 / LP5816 & \ac{LED} drivers, status \acp{LED} and \ac{LCD} backlight & \ac{I2C}4 (0x2D/0x2C) \\
        OPT3004 & Ambient light sensor, automatic backlight & \ac{I2C}4 (0x44) \\
        \ac{LCD} (240$\times$320) & Display & \ac{SPI}4 (hardware, \ac{DMA}) \\
    \end{tabular}
    \caption{Key external \acp{IC} and their bus interfaces}
\end{table}

\ac{ADC}1--5 each own a \ac{DMA} channel: \ac{ADC}1 carries \texttt{I\_OUT}
and the control loop, \ac{ADC}2 \texttt{V\_TERM}, \ac{ADC}3 \texttt{V\_IN},
\ac{ADC}4 \texttt{V\_OUT}/\texttt{V\_HV} (switched per mode), and \ac{ADC}5
a 12-channel scan (rails, \acp{NTC}, VBAT, VREFINT) plus an injected group
for \texttt{I\_ISO}/\texttt{I\_BAT}, oversampled 256$\times$ for an
effective 16-bit result. \texttt{adc\_configure\_mode()} re-points triggers
and channels per operating mode.

Two \ac{DAC} channels are used: \texttt{DAC1\_CH1} drives an analog
reference voltage, \texttt{DAC1\_CH2} drives the \qty{1}{\ampere} reference
for the milliohm mode as a \qty{1}{\hertz} \qty{50}{\percent} square wave
(TIM6), so the reading can be gated to the on-phase.

All \ac{I2C}4 slaves share one bus; transactions are queued and
\ac{DMA}/event-driven, dropping with an error code rather than blocking
when the queue is full. \ac{SPI}3 talks to the external \ac{ADC}; \ac{SPI}4
is a real hardware \ac{SPI} peripheral driving the \ac{LCD}, with bulk
transfers \ac{DMA}-driven -- only the chip-select and data/command lines
are toggled by \ac{GPIO}, which is the usual wiring for an \ac{SPI} \ac{LCD},
not bit-banging. \ac{HRTIM}
drives three channels -- primary buck-boost, high-voltage flyback, and the
secondary winding -- with hardware dead-time insertion on the primary and
secondary legs (see \autoref{sec:fmea}).

Interrupts run in three fixed priority tiers (\texttt{irq\_priority.h}):
the control loop's \ac{DMA} channel alone at the top, the remaining
\ac{ADC}/\ac{SPI}3/external-\ac{ADC}-ready interrupts next, and \ac{I2C},
\ac{LCD} \ac{SPI} and the state-machine timer last. Because the control
loop can preempt the state-machine step, \texttt{ctrl\_main\_start\_ctrl()}
and \texttt{ctrl\_main\_stop\_control()} order their writes to \texttt{mode}
so a preempting tick always observes either \texttt{CTRL\_MODE\_OFF} or a
fully configured mode, never a half-built one.

\subsection{Persistent configuration and calibration}

\texttt{config\_store.c} keeps a versioned layout on the \ac{I2C}
\ac{EEPROM}: an 8-byte header, a 16-byte hardware-data block (serial
number) and a 192-byte calibration block, each with its own \ac{CRC}32.
\texttt{calibration.c} exposes six calibratable channels (\texttt{V\_TERM},
\texttt{V\_SENS}, \texttt{V\_OUT}, \texttt{V\_HV}, \texttt{I\_OUT},
\texttt{I\_ISO}), reachable from Settings > Calibration or the \ac{CLI}
(\texttt{zeroCal}/\texttt{gainCal}/\texttt{setOffset}/\texttt{setGain}), and
persists each step immediately. \autoref{sec:fmea} covers what happens on a
corrupted block.

\subsection{Protection}

\texttt{protection.c} polls the \ac{OVP}/\ac{OCP} fault-latch pins and four
external \ac{NTC} temperatures once per state-machine tick, classifying
each temperature into OK/WARNING/ERROR with hysteresis. An ERROR of any
kind forces the converter output off from \texttt{statemachine\_step()}
regardless of mode. The mechanism is summarised here; the full
failure-by-failure treatment, including where this protection does and
does not currently reach, is in \autoref{sec:fmea}.

\subsection{User interface}

A 240$\times$320 \ac{TFT} \ac{LCD} is driven through a custom driver and
$\mu$\ac{GUI}. Backlight brightness (LP5816) fades toward a target derived
from the OPT3004 ambient-light sensor, in lux. A rotary encoder plus
three buttons (ESC, OUT, OK) are debounced in \texttt{input.c}. Menu icons
come from bitmaps in the QSPI flash's icon store when available, falling
back to procedural vector glyphs otherwise. Status \acp{LED} and the
backlight are managed by \texttt{ui\_ctrl.c}.

\subsection{Command-line interface}

\texttt{cli.c} implements a polled command parser over USART2, used for
calibration, diagnostics and parameter read/write during bench work. It
includes a keyboard-driven \texttt{control} mode that maps terminal keys to
the encoder/buttons, so the \ac{UI} can be exercised without the display
board attached.

\subsection{Build and verification}

The firmware is built with STM32CubeIDE (managed build, \texttt{-Os}).
There is no unit-test harness and no host build. The available automated
check is \texttt{sw/bat\_source/tools/check.sh}: it syntax-compiles every
\texttt{Core/Src/*.c} file with \texttt{-Wall -Wextra} using the same flags
as the project file, and diffs the resulting warnings against a checked-in
baseline, failing on anything new. It does not link, so section sizes and
stack headroom still require a full CubeIDE build.

\subsection{Open items}

Safety-relevant open items are tracked as failure modes in
\autoref{sec:fmea}. Remaining open items with no safety impact --
PID bench tuning, the calibration screen not restoring \ac{ADC} routing on
exit, and data logging to the QSPI flash not yet being wired up -- are
tracked in \texttt{sw/architecture.md}, \S10.
